a group homomorphism from @@M0@@ to @@M1@@; the map which assigns to a series its @@M2@@-st coefficient, multiplied by @@M3@@, is an additive homomorphism from @@M4@@ to @@M5@@; the morphism @@M6@@ will be the composite of the two. It is clear that @@M7@@ by definition (7.1.1).

The second of formulas (7.5.1) therefore expresses the commutativity of the diagram written above. If we suppose that @@M8@@ is a @@M9@@-ring, the left square is commutative, and it suffices to verify the commutativity of the right one. Now this is a diagram of homomorphisms of functors on @@M10@@, with source @@M11@@; it therefore suffices, by 1.10, to verify commutativity for @@M12@@; we then have:
@@M13@@
@@M14@@, which gives the result.

To prove the last one, we first observe that the diagram:
\sgaStableTarget{sga6.e05.eq.7-5-3}@@M15@@
is commutative. As above, it indeed suffices to see this for @@M16@@, and we have:
@@M17@@
@@M18@@

Consequently, @@M19@@.

Moreover, @@M20@@.

